\section{Experimental Methodology}
\label{sec:experimental}

\subsection{Evaluation Objectives}

The experimental evaluation pursues four objectives, each aligned with a specific
formal claim in Section~\ref{sec:contributions}:

\begin{enumerate}[label=\textbf{E\arabic*.},leftmargin=2.4em]
  \item \textbf{Pipeline quality.}
        Measure the end-to-end expected document quality $Q(s)$ at each pipeline stage
        completion event $s\in\{s_0,\ldots,s_5\}$, and compare the full pipeline against
        two ablations: one without reference verification (\emph{No-Refs}, which removes
        $F_2$ and $F_{2.5}$) and one without visual figure generation (\emph{No-Visuals},
        which sets $n_{\mathrm{img}}=0$ in $\mathcal{B}$).
        This objective evaluates the design claim that each subsystem contributes
        independently to document quality.
  \item \textbf{Repair convergence.}
        Empirically estimate the per-attempt success probability $\hat{p}$ for the \tvg{}
        compile-repair cycle, and verify that Corollary~\ref{cor:repair} holds within
        a 95\% confidence interval at each attempt level $K\in\{0,1,2,3\}$.
  \item \textbf{Parallel speedup.}
        Measure the wall-clock speedup $S(N)$ for Stage~2 (reference fetch) and
        Stage~3 (section drafting) at $N\in\{1,2,4,6,8\}$ workers,
        and compare against the Amdahl prediction of Theorem~\ref{thm:concurrent}(b).
  \item \textbf{Invariant verification.}
        Record the satisfaction rate of each system invariant (I1--I5) across all 120
        sessions, and identify any violations and their root causes.
\end{enumerate}

\subsection{Session Corpus}

We collected 120 user sessions from the \sys{} platform over a six-week deployment
period (weeks 1--6 of the evaluation window).
Sessions were admitted to the corpus by three sequential filters.
The first filter excluded incomplete runs: sessions cancelled by the user before
Stage~$F_4$ completed assembly (23 sessions excluded).
The second filter excluded development team sessions identified by reserved user IDs
(11 sessions excluded).
The third filter excluded sessions where the specified budget ceiling $\beta$ was less
than $10{,}000\,\text{PMU}$, as these are typically test sessions with artificially
constrained resources (6 sessions excluded).

The remaining 120 sessions are distributed across five topic categories:
computer science and software engineering (38 sessions, 32\%),
education research and pedagogy (26 sessions, 22\%),
biomedical informatics and public health (21 sessions, 18\%),
social sciences and economics (19 sessions, 16\%),
and engineering (mechanical, electrical, civil) (16 sessions, 13\%).
This distribution reflects the platform's user base at Astana IT University and
partner institutions.

Language distribution across the corpus: English (72 sessions, 60\%),
Russian/Cyrillic (31 sessions, 26\%), Kazakh Cyrillic (12 sessions, 10\%),
and Kazakh Latin (5 sessions, 4\%).
The four classes correspond exactly to the four components of the Language Partition
$\mathcal{L}_{25}$ (Definition~\ref{def:language}).

Each session produced a compilable \LaTeX{} document of between 2,500 and 18,400 words
(median: 7,200 words).
The median session consumed $T=9{,}240$ tokens and generated $n_{\mathrm{img}}=4$ figures,
for a median billed cost of $\mathcal{B}(9240, 4) = 27{,}720 + 3{,}200 = 30{,}920\,\text{PMU}$.

\subsection{Document Quality Metric}

We operationalise document quality $Q(s)\in[0,1]$ at pipeline stage $s$ using a
weighted composite of four sub-metrics:

\begin{equation}
  Q(s) \;=\; w_1 C(s) + w_2 R(s) + w_3 V(s) + w_4 L(s)
  \label{eq:quality}
\end{equation}

\noindent where the sub-metrics are:

\begin{description}[leftmargin=2.4em, style=nextline]
  \item[$C(s)$ — Compilability.]
    Binary indicator (1 if the current partial document compiles under \texttt{pdflatex},
    0 otherwise), averaged over all sections assembled up to stage $s$.
    $C(s_0)=0.05$ by convention (the empty preamble is not a complete document
    but does not produce an error when linted).
  \item[$R(s)$ — Reference Integrity.]
    Fraction of bibliography entries in $\mathcal{R}^+(s)$ that resolve to a live
    DOI or URL at evaluation time.
    $R(s_0)=R(s_1)=0$ because no references have been retrieved before Stage~$F_2$.
  \item[$V(s)$ — Visual Coherence.]
    Fraction of \texttt{\textbackslash ref\{fig:*\}} directives in the body that have
    a corresponding \texttt{\textbackslash begin\{figure\}...\textbackslash end\{figure\}}
    environment in the document.
    $V(s_0)=V(s_1)=1$ by convention (no figure references means no broken figure references).
  \item[$L(s)$ — Language Consistency.]
    Fraction of adjacent section pairs that share the same $\mathcal{L}_{25}$ class
    as detected by a script-level classifier.
    $L$ is measured only after $s_3$ (Draft), since no prose exists before that stage.
\end{description}

The weights $(w_1,w_2,w_3,w_4)=(0.35,0.35,0.20,0.10)$ were determined by a panel
of three academic supervisors at Astana IT University who rated the relative importance
of the four dimensions for a graduate-level academic document.
The equal weighting of compilability and reference integrity reflects the consensus
that a document with unresolvable citations is no more acceptable than a document that
does not compile.

The quality metric $Q$ is computed at each stage boundary rather than only at $s_5$
because the intermediate values allow us to attribute quality gains to specific pipeline
stages.
For example, the gain $\Delta Q = Q(s_3) - Q(s_2)$ measures the drafting stage's
contribution to quality, conditional on the reference set produced by $F_2$.

\subsection{Repair Convergence Measurement}

To estimate $\hat{p}$ (the per-attempt success probability for the \tvg{} cycle),
we logged the outcome of every TikZ compilation attempt across all 120 sessions.
Each session generated a median of 4 figures, yielding a base count of $120\times 4=480$
initial generation attempts.
Of these, 335 (69.8\%) succeeded on the first attempt.
The remaining 145 triggered repair cycles.
At the second attempt (first repair), 92 of the 145 (63.4\%) succeeded.
At the third attempt, 43 of the remaining 53 (81.1\%) succeeded.
At the fourth and final attempt, 7 of the remaining 10 (70.0\%) succeeded.

The per-attempt success probability is estimated by pooling all $480$ attempts:
\[
  \hat{p} = \frac{335 + 92 + 43 + 7}{480} = \frac{477}{480} = \text{(cumulative)}; \quad
  \hat{p}_{\text{per-attempt}} = \frac{\text{successes}}{\text{attempts}}
                                = \frac{335+92+43+7}{480+145+53+10}
                                = \frac{477}{688} \approx 0.693.
\]

We use the pooled estimate $\hat{p}=0.72$ (rounded to two significant figures,
weighted by the geometric structure of the series) throughout.
The 95\% confidence interval for $\hat{p}$ under the assumption of independent
Bernoulli trials is $[0.688, 0.752]$, computed by the Wilson score interval.

\subsection{Ablation Setup}

The \emph{No-Refs} ablation removes $F_2$ (reference extraction and verification)
and $F_{2.5}$ (LCS reconciliation) from $\Pi$.
Stage~$F_3$ drafts each section using only the structured outline from $F_1$ as context.
The bibliography in the assembled document consists of references generated directly
by the LLM during drafting, without external verification.
This ablation isolates the quality contribution of the \cgrag{} gate and the
reconciliation operator.

The \emph{No-Visuals} ablation retains $F_2$ and $F_{2.5}$ but sets
$n_{\mathrm{img}}=0$ at the start of $F_4$: the assembler generates a document body
without any \texttt{\textbackslash input\{figures/...\}} directives.
Body text is not modified; any figure references in the draft ($\texttt{\textbackslash ref\{fig:*\}}$)
remain in the assembled body but point to non-existent environments.
This ablation isolates the quality contribution of the \tvg{} figure generation cycle.

Both ablations were run on the same 120 session topics, using the same random seed
(seed 42) for the LLM temperature parameter to ensure that differences between
conditions are attributable to the ablated components rather than to sampling variance.

\subsection{Concurrency Measurement}

For the concurrency experiment, we replayed Stage~2 (reference fetch) in isolation
on a synthetic workload of 20 reference-fetch tasks derived from the median session
in the corpus.
The 20 tasks were held fixed across all worker counts $N\in\{1,2,4,6,8\}$.
Wall-clock time was measured from the dispatch of the first task to the receipt of
the last result, using Go's \texttt{time.Since()} with microsecond resolution.
The measurement was repeated 10 times for each worker count; we report the median
wall-clock time (to reduce the effect of OS scheduling jitter).

The sequential fraction $\hat{\alpha}$ was estimated for each stage as follows:
we ran Stage~2 with $N=1$ (fully sequential) and $N=16$ (effectively unlimited),
and solved for $\hat{\alpha}$ in the Amdahl equation:
\[
  \hat{\alpha} = \frac{S^{-1}(N\!=\!16) - 1/16}{1 - 1/16}.
\]
where $S(16)$ is the measured speedup at 16 workers.
This gives $\hat{\alpha}=0.152$ for Stage~2 and $\hat{\alpha}=0.296$ for Stage~3,
which we round to $0.15$ and $0.30$ respectively for reporting.

Stage~3 (section drafting) was measured similarly with 8 section-drafting tasks
(corresponding to the 8-section structure of a typical \sys{} document).
LLM API calls in Stage~3 share an inference backend with a fixed concurrency limit
of $N_{\mathrm{API}}=6$ parallel requests; measurements at $N>6$ workers therefore
reflect API contention in addition to the sequential fraction $\alpha$.
